\documentclass[10pt,a4paper]{amsart}
\usepackage[margin=19mm]{geometry}
\usepackage{amsmath,amssymb,mathtools}
\usepackage[scr=rsfs,cal=euler]{mathalfa}
\usepackage{xfrac}
\usepackage[unicode,hidelinks,bookmarks=false]{hyperref}
\newcommand{\ie}{i.e.\ }
\newcommand{\cf}{cf.\ }
\newcommand{\resp}{respectively\ }
\newcommand{\Subsection}{\subsection}
\providecommand{\nobreakdash}{\nobreak\hbox{-}}
\setlength{\emergencystretch}{2em}
\multlinegap0pt
\usepackage{tikz}
\usepackage{subfigure}
\usepackage{ytableau}
\usepackage{array}
\usepackage{mathtools}
\setlength{\textheight}{220mm}
\setlength{\textwidth}{155mm}
\setlength{\oddsidemargin}{2.1mm}
\setlength{\evensidemargin}{2.1mm}
\setlength{\topmargin}{0mm}
\renewcommand{\baselinestretch}{1.2}
\newtheorem{theorem}{Theorem}[section]
\newtheorem{lemma}[theorem]{Lemma}
\newtheorem{corollary}[theorem]{Corollary}
\newtheorem{proposition}[theorem]{Proposition}
\newtheorem{conjecture}[theorem]{Conjecture}
\newtheorem{definition}{Definition}
\newtheorem{example}[]{Example}
\newtheorem{remark}[]{Remark}
\newtheorem{application}[]{Application}
\renewcommand{\Re}{\operatorname{Re}}
\renewcommand{\Im}{\operatorname{Im}}
\newcommand{\df}{\dfrac}
\newcommand{\Arg}{\operatorname{Arg}}
\newcommand{\Res}{\operatorname{Res}}
\renewcommand{\a}{\alpha}
\renewcommand{\b}{\beta}
\newcommand{\e}{\epsilon}
\renewcommand{\d}{\delta}
\newcommand{\g}{\gamma}
\newcommand{\G}{\Gamma}
\renewcommand{\l}{\lambda}
\renewcommand{\o}{\omega}
\renewcommand{\r}{\rho}
\renewcommand{\t}{\tau}
\newcommand{\z}{\zeta}
\newcommand{\qu}{{\underline{q}}}
\numberwithin{equation}{section}
\begin{document}
\title[Asymptotic Formula for Multipartitions]{Asymptotic Formula for Multipartitions}
\author{JAYANTA BARMAN JAYANTA BARMAN; Jayanta Barman; Kamalakshya Mahatab}
\date{}
\maketitle
\noindent\emph{Source and licence.} Asymptotic Formula for Multipartitions. Version arXiv:2607.09159v1 (CC BY 4.0). This material is distributed under the Creative Commons Attribution 4.0 International licence, http://creativecommons.org/licenses/by/4.0/, as recorded in the arXiv OAI licence metadata for this exact version and shown on the abstract page https://arxiv.org/abs/2607.09159v1. Full rights notice: licenses/arxiv-2607.09159-CC-BY-4.0.txt.\par\smallskip
\noindent\textit{Editorial scope.} Counted unit: the original \texttt{theorem} environment \texttt{theorem-1.2} at source lines 143--154 together with its complete proof at lines 483--526; the delivered document keeps source lines 75--537 so that every referenced label and every citation resolves inside it. Native editable source is the paper's own concatenated TeX; the AMS wrapper is editorial formatting only. No publisher class, upstream script or unrelated artwork is redistributed.
\par\smallskip\noindent\textit{Reference note.} This article ships no compiled bibliography (biblatex); the entries delivered here are composed from the author's own BibTeX (.bib) fields, with no field invented and no entry dropped.
	\section{Introduction}\label{section1} 
	The asymptotic theory of partition functions has its origin in the celebrated work of Hardy and Ramanujan \cite{hardy1918asymptotic}, who introduced the circle method and obtained an asymptotic formula for the ordinary partition function $p(N)$. 
	%	\begin{equation*}
		%		p(N)\sim \frac{1}{4N\sqrt{3}}\exp\left(\frac{2\pi}{\sqrt{6}}\sqrt{N}\right)
		%		\quad \text{as } N\to\infty.
		%	\end{equation*}
	%	Later, Rademacher \cite{rademacher1938partition} refined this result into an exact convergent series:
	%	\begin{align*}
		%		p(N)= \frac{1}{4N\sqrt{3}}\exp\left(\frac{2\pi}{\sqrt{6}}\sqrt{N}\right)\left(1+O\left(N^{-\frac{1}{2}}\right)\right).
		%	\end{align*}
	Although unrestricted partitions form a central object of study, imposing constraints on partitions opens many deep and interesting questions in both arithmetic and combinatorics.
	One classical example is that of $t$-multipartitions: a $t$-multipartition of a positive integer $N$ is a sequence $\left(\mu^{(1)},\mu^{(2)},\ldots,\mu^{(t)}\right)$, where each $\mu^{(j)}$ is an integer partition (possibly empty) for all $j$ such that $\sum_{j=1}^{t}\left|\mu^{(j)}\right|=N$. Here, $\left|\mu^{(j)}\right|$ denotes the sum of the parts of the partition $\mu^{(j)}$. Let $p_{t}(N)$ denote the number of $t$-multipartitions of $N$. Further, the number of $t$-multipartitions of $N$ coincides with the number of $t$-colored partitions of $N$. Equivalently, a $t$-colored partition of $N$ is a partition in which each part may be assigned one of $t$ available colors, with the order of the colors being immaterial. For instance, the $2$-colored partitions of $3$ are $3_1, 3_2, 2_1+1_1, 2_1+1_2, 2_2+1_1,2_2+1_2,1_1+1_1+1_1, 1_2+1_1+1_1, 1_2+1_2+1_1, 1_2+1_2+1_2.$ 
	Since these two notions are equivalent from an enumerative point of view, we may also interpret $p_t(N)$ as the number of $t$-colored partitions of $N$. Various properties of $p_t(N)$, including log-concavity, arithmetic properties, and multiplicative properties, have been studied in \cite{bringmann, bringmann1, chern}. The generating function for $t$-multipartitions, together with their congruence and other arithmetic properties, has been investigated extensively in \cite{andrews, atkin, chen2014}. 
	
	James and Kerber \cite[Corollary 4.4.4]{MR644144} showed that $p_t(N)$ can be expressed as a sum of products of partition functions, given by
	\begin{equation*}
		\sum_{\substack{N=n_{1}+n_{2}+\cdots+n_{t}\\ n_{j}\ge 0}} p(n_{1})p(n_{2})\cdots p(n_{t}).
	\end{equation*}
	This is equal to the number of irreducible representations of the wreath product $G \wr S_N$, where the group $G$ has $t$ conjugacy classes and $S_N$ is the symmetric group. There are also applications of $p_t(N)$ in the representation theory of Lie algebras \cite{bouwknegt2002,fayers2006weights}. In \cite{bouwknegt2002}, multipartitions play an important role in the study of Durfee systems and their existence. Furthermore, \cite{fayers2006weights} shows that the irreducible representations of the Ariki--Koike algebra are naturally indexed by $t$-multipartitions of $N$. Multipartitions are also important in algebraic geometry. For a smooth projective surface $S$, let $S^{[N]}$ be the Hilbert scheme of $N$ points, which is a smooth projective variety of dimension $2N.$ Let $\chi(S)$ and $\chi(S^{[N]})$ be the topological Euler characteristic of $S$ and $S^{[N]}$, respectively. Then (see \cite[equation $(2)$, Theorem 0.1]{gottsche}),
	\begin{align*}
		\chi\left(S^{[N]}\right)=p_{\chi(S)}(N).
	\end{align*}
	For applications of multipartitions to gauge theory and random partitions, see \cite{nekrasov}. 
	%	Formulas of this type naturally lead to questions about the asymptotic behavior of multipartition functions. Although the circle method is one of the principal tools in analytic number theory, its implementation is often technically complicated. Simpler approaches to partition asymptotics are therefore valuable for both their mathematical significance and their expository appeal. 
	
	
	In \cite{murty2013partition}, Murty applied the Laplace saddle point method to derive an asymptotic formula for $p_t(N)$ for fixed $t$ without providing any error term. This result also follows from Meinardus's formula \cite{meinardus1953}, where special functions play an essential role. In this article, we extend this analysis to the range $t \ll N^{1-\epsilon}$ for any $\epsilon>0$, instead of a fixed $t$, using the saddle point method. This method also has been successfully used to obtain asymptotic formulas for partition functions and related combinatorial sequences \cite{barman2025lower, barman2026asymptotic, tyler2026asymptotics}. %We present a simplified proof that includes an explicit error term, thereby filling a gap in the existing literature, where such effective results were previously unavailable. 
	
	%  The generating function for $p(N)$ is given by
	%  \begin{equation*}
		%  F(q)=\prod_{n=1}^{\infty}(1-q^{n})^{-1}=\sum_{N=0}^{\infty}p(N)q^{N}=q^{\frac{1}{24}}\frac{1}{\eta(z)}.
		%  \end{equation*}
	%  By Cauchy's integral formula, the $N$-th coefficient of the series is given by
	%  \begin{equation*}
		%  p(N)=\frac{1}{2\pi i}\int_{C}\frac{F(q)}{q^{N+1}}dq 
		% \end{equation*}
	%  where $C$ is a simple, positively oriented loop around the origin, lying entirely on the unit circle. In this article, we use the saddle point method to evaluate the integral
	% \begin{equation}
		% p(N)=\int_{-1/2}^{1/2}\exp\left(-2\pi iz\left(N-\frac{1}{24}\right)\right)h(z)dx,   
		% \end{equation}
	% where $h(z)=\frac{1}{\eta(z)}$.
	Now to state our main theorem and give an idea of its proof, we recall some identities involving $p_t(N)$. 
	The generating function for $p_{t}(N)$ is given by
	\begin{equation*}
		G(q)=\prod_{N=1}^\infty \left(1-q^{N}\right)^{-t}=\sum_{N=0}^{\infty}p_t(N)q^{N}=\left(\prod_{N=1}^\infty \frac{1}{\left(1-q^{N}\right)}\right)^{t}=\frac{q^{\frac{t}{24}}}{\eta(z)^{t}}, 
	\end{equation*}
	where $q=\exp(2\pi iz)$, $z=x+iy$ and $y>0$. The Dedekind eta function $\eta(z)$ is given by 
	\begin{equation*}
		\eta(z)=\exp\left(\frac{\pi iz}{12}\right)\prod_{n=1}^{\infty}(1-\exp(2\pi inz)).  
	\end{equation*}
	We will use the functions $\mu_m$, $m\ge 1$, from \cite{tyler2026asymptotics} throughout this article:
	\begin{equation}\label{eq-muk}
		\mu_{m}(z)=-\frac{z^{m+1}}{2\pi i} \left(\frac{d}{dz}\right)^{m} \log\eta(z).
	\end{equation}
	By the Cauchy integral formula, the $N$-th coefficient of the series is given by
	\begin{equation*}
		p_{t}(N)=\frac{1}{2\pi i}\int_{\gamma}\frac{G(q)}{q^{N+1}}dq, 
	\end{equation*}
	where $\gamma$ is a simple positively-oriented loop around the origin, located entirely in the unit circle. Let $\gamma$ be the contour defined by $q=\exp(2\pi i z)$. The condition $|q|<1$ holds if and only if $y>0$. For a fixed value of $y$, as $x$ varies over any interval of length $1$, the variable $q$ moves along a full circle of radius $e^{-2\pi y}$. Therefore, we can write $p_t(N)$ as
	\begin{align}\label{eq-integral}
		p_{t}(N)&=\int_{-1/2}^{1/2} \exp\left( -2\pi izM\right)a_{t}(z)dx,
	\end{align}
	where 
	\begin{align*}%\label{eq-atz}
		a_{t}(z)=\frac{1}{\eta(z)^{t}}\quad\text{and}\quad M=N-\frac{t}{24}.
	\end{align*}	
	
	We now state the main theorem of this article.

	\begin{theorem}\label{theorem-1.2}
		Let $t$ be a positive integer $<N$.
		\newline
		(i) Then there exists a unique solution $y>0$, such that
		\begin{align}\label{eq-1.1}
			y^{2}M+t\mu_{1}(iy)=0.
		\end{align}
		(ii) Corresponding to this value of $y$, the function $p_t(N)$ satisfies the following asymptotic approximation:
		\begin{align*}
			p_t(N)=\frac{y^{\frac{3}{2}}\exp(2\pi My)a_{t}(iy)}{\sqrt{t\mu_{2}(iy)}}\left(1+O\left(\frac{y}{t}\right)\right).    
		\end{align*}
	\end{theorem}

	Since $y$ is not given explicitly, we solve for $y$ in terms of $N$ and $t$ using Theorem~\ref{theorem-1.2}$(i)$ and substitute it in part~$(ii)$. This gives the following asymptotic formula for $p_t(N)$.
	\begin{theorem}\label{theorem-1.1}
		For any $\epsilon > 0$, let $t \ll N^{1 - \epsilon}$ be a positive integer. Then
		\begin{align*}
			p_t(N)=&\left(\frac{t}{24}\right)^{\frac{t+1}{4}}\frac{\exp\left(\frac{2\pi}{\sqrt{6}}\sqrt{t}\sqrt{N-\frac{t}{24}}\right)\exp\left(O\left(\frac{t^{\frac{3}{2}}}{\sqrt{N}}\right)\right)}{\sqrt{2}\left(N-\frac{t}{24}\right)^{\frac{t+3}{4}}}\left(1 + O\left(\frac{1}{\sqrt{Nt}}\right)\right). 
		\end{align*}
	\end{theorem}
	
	%	\subsection{\texorpdfstring{Applications of Theorem \ref{theorem-1.1}}{ Applications of Theorem 1.2}}In this subsection, we explore several applications of Theorem \ref{theorem-1.1}.
	In Theorem \ref{theorem-1.1}, if we set $t = 1$, we obtain Rademacher's asymptotic formula \cite{rademacher1938partition} for $p(N)$.
	\begin{corollary}\label{corollary-1.3}
		Let $N$ be a large positive integer. Then  
		\begin{equation*}
			p(N)=\frac{1}{4\sqrt{3}N}\exp\left(\frac{2\pi}{\sqrt{6}}\sqrt{N}\right)\left(1+O\left(N^{-\frac{1}{2}}\right)\right).
		\end{equation*}
	\end{corollary}
	In Theorem \ref{theorem-1.1}, if we take $t$ to be any fixed positive integer, then we obtain the following corollary, including the error term. This result coincides with Murty's result \cite{murty2013partition} for $p_t(N)$.
	\begin{corollary}[{\cite[Theorem 4]{murty2013partition}}]\label{corollary-1.4}
		Let $t$ be any fixed
		positive integer. Then
		\begin{equation*}
			p_{t}(N)=\left(\frac{t}{24}\right)^{\frac{t+1}{4}}\frac{\exp\left(\frac{2\pi}{\sqrt{6}}\sqrt{Nt}\right)}{\sqrt{2}N^{\frac{t+3}{4}}}\left(1+O\left(N^{-\frac{1}{2}}\right)\right).
		\end{equation*}  
	\end{corollary}
	We may compare Corollary~\ref{corollary-1.4} with Meinardus's formula \cite{meinardus1953}:
	\begin{align*}
		p_t(N)=\frac{\exp\left(D'(0)\right)\left(t\Gamma(2)\zeta(2)\right)^{\frac{1-2D(0)}{4}}\exp\left(2\sqrt{Nt\Gamma(2)\zeta(2)}\right)}{\sqrt{4\pi}\,N^{\frac{3-2D(0)}{4}}} \left(1+O\left(N^{-\frac{1}{2}}\right)\right), 
	\end{align*}
	where $D(s)=t\zeta(s)$, and $\zeta(s)$ is the Riemann zeta function. This gives an alternative proof of special values of $\zeta(s)$: 
	$$
	\zeta(0)=-\frac{1}{2}, \qquad
	\zeta'(0)=-\frac{1}{2}\log(2\pi), \qquad
	\zeta(2)=\frac{\pi^2}{6}.
	$$ 
	
	%	Let $p_{t}(N, l)$ denote the number of $t$-multipartitions of $N$ into parts of size at most $l$. Then the number of $t$-multipartitions for which the largest part has size 
	%	\begin{equation*}
		%		l \ge \frac{\sqrt{6}}{2\pi} \sqrt{\frac{N(1+\delta)}{t}} \left(\log \frac{N}{t}(1+\delta)\right) \left(1 + \frac{1}{A}\right)
		%	\end{equation*}
	%	is given by $p_{t}(N) - p_{t}(N, l)$. By the proof of Theorem 3.9 in \cite{dong2023almost}, the number of zeros in the character table of $G \wr S_{N}$ is at least
	%	\begin{equation}\label{eq-Zeros}
		%		p_{t}(N)\left(1 - O\left(\frac{\log\left(\frac{N}{t}(1+\delta)\right)}{\left(\frac{N}{t}(1-\delta)\right)^{\frac{1}{2A}}}\right)\right)\left(p_{t}(N) - p_{t}(N, l)\right),
		%	\end{equation}
	%	where $G$ is a finite group with $t$ conjugacy classes, $\delta > 0$, and $A \le \frac{\log |\lambda_{i}|}{\log\log |\lambda_{i}|}$ for each $i$.
	%	Therefore, using Theorem \ref{theorem-1.1}, if one can estimate the number of $t$-multipartitions of $N$ that contain at least one part greater than
	%	\begin{equation*}
		%		\frac{\sqrt{6}}{2\pi} \sqrt{\frac{N(1+\delta)}{t}} \left(\log \frac{N}{t}(1+\delta)\right) \left(1 + \frac{1}{A}\right).
		%	\end{equation*}
	%	Then one can also estimate a lower bound for the number of zeros in the character table of the wreath product $G \wr S_N$ using (\ref{eq-Zeros}).
	
	
	\subsection{\texorpdfstring{Sketch of the proof of Theorem \ref{theorem-1.2}}{}} 
	We use the integral in (\ref{eq-integral}) as a basis for our saddle point analysis. 
	Note that  $y=\Im (z)$ only appears on the right-hand side of (\ref{eq-integral}).
	Hence, we select $y$ appropriately as in (\ref{eq-1.1}) to obtain an asymptotic formula for $p_t(N)$. 
	The Taylor series of $a_t(z)=\frac{1}{\eta(z)^{t}}$ has radius of convergence $<y$, so we divide the integral in (\ref{eq-integral}) into two separate parts:
	\begin{align*}
		|x| < \frac{y}{3} 
		\qquad \text{and} \qquad 
		\frac{y}{3} \le |x| \le \frac{1}{2}.
	\end{align*}
	The integral in (\ref{eq-integral}) may be decomposed as
	\begin{equation}
		\begin{aligned}\label{eq-ptN1}
			p_{t}(N)&=\exp(2\pi My)a_{t}(iy)\int_{-y/3}^{y/3}\exp\left(-2\pi i Mx+2\pi i\frac{1}{2\pi i}\log\frac{a_{t}(z)}{a_{t}(iy)}\right)dx\\
			&+\exp(2\pi My)a_{t}(iy)\int_{\frac{y}{3}\le |x|\le \frac{1}{2}}\exp\left(-2\pi i Mx\right)\frac{a_{t}(z)}{a_{t}(iy)}dx.
		\end{aligned}
	\end{equation}   
	The Taylor expansion of $\log a_t(z)$ around $x=0$ is given by
	\begin{align}\label{eq-tyler1}
		\log a_{t}(z)&=\left(\log a_{t}(iy)+x\frac{d}{dz} \log a_{t}(iy)+\frac{x^{2}}{2!}\frac{d^{2}}{dz^{2}}\log a_{t}(iy)+\cdots\right).
		%		\\
		%		\log \frac{a_{t}(z)}{a_{t}(iy)}&=\left(x\frac{d}{dz} \log a_{t}(z)+\frac{x^{2}}{2!}\frac{d^{2}}{dz^{2}}\log a_{t}(z)+\cdots\right)_{z=iy}.
	\end{align}
	From the definition of $\mu_m$, we have
	\begin{equation*}
		\left(\frac{d}{dz}\right)^{m}\log a_{t}(z)=t\frac{2\pi i}{z^{m+1}}\mu_{m}(z).
	\end{equation*}
	So, (\ref{eq-tyler1}) simplifies to
	\begin{align}\label{eq-taylor}
		\frac{1}{2\pi i}\log \frac{a_{t}(z)}{a_{t}(iy)}
		&=t\left(x\frac{\mu_{1}(iy)}{(iy)^{2}}+\frac{x^{2}}{2}\frac{\mu_{2}(iy)}{(iy)^{3}}+\frac{x^{3}}{6}\frac{\mu_{3}(iy)}{(iy)^{4}}+\frac{x^{4}}{24}\frac{\mu_{4}(x^{\prime}+iy)}{(x^{\prime}+iy)^{5}}\right),
	\end{align}
	for some $z^{\prime}=x^{\prime}+iy$, where $x^{\prime}$ is between $0$ and $x$. In the region $|x|<\frac{y}{3}$, we evaluate the integral using the Taylor expansion described above. However, this expansion is not valid in the region $\frac{y}{3} \le |x| \le \frac{1}{2}$. Instead, we establish Lemma \ref{proposition-error}, where we derive an upper bound. Combining these results, Proposition \ref{theorem-main} shows that $|x|<\frac{y}{3}$ contributes to the main term, and $\frac{y}{3}\le |x|\le \frac{1}{2}$ contributes to the error in (\ref{eq-ptN1}).
	
	
	%	\begin{lemma}
		%		Let $t \ge 1$ be any positive integer. Then $p_t(N) < p_{t+1}(N)$ for all $N \ge 1$.
		%	\end{lemma}
	%	\begin{proof}
		%		Let $\mu=\left(\mu^{(1)},\mu^{(2)},\ldots,\mu^{(t)}\right)$ be an arbitrary $t$-multipartition of $N$. If we take $\mu^{(t+1)}=(0)$, then $\mu'=\left(\mu^{(1)},\mu^{(2)},\ldots,\mu^{(t)},\mu^{(t+1)}\right)$ is a $(t+1)$-multipartition of $N$. Since we take $\mu$ ia an arbitrary $t$-multipartition of $N$. So, $p_t(N)\le p_{t+1}(N)$. Since $\sum_{j=1}^{t}\left|\mu^{(j)}\right|=N$, atlaest one of the components, $|\mu^{(k)}|$, $1\le k\le t$ is non zero. We make another $(t+1)$-multipartition by taking $\mu^{(t+1)}=\mu^{(k)}$ and then $\mu^{(k)}=(0)$. So, we got a new $(t+1)$-multipartition $\mu''=\left(\mu^{(1)},\ldots,\mu^{(k-1)},(0),\mu^{(k+1)},\ldots\mu^{(t)},\mu^{(k)}\right)$. So, $p_t(N)<p_{t+1}(N)$.   
		%	\end{proof}
	\section{Acknowledgments}
	J. Barman is deeply thankful to the University Grants Commission (UGC), India, for their invaluable support through the Fellowship Programme. K. Mahatab is supported by the ARG-MATRICS Programme (grant no. ANRF/ARGM/2025/002540/MTR).
	\section{Preliminaries}
	In this section, we establish several preliminary results that are required to prove the main results of this paper. 	
	%	We use the following lemma to bound $\eta(z)$ when $y \ge \frac{\sqrt{3}}{2}$ is large.
	%	\begin{lemma}[{\cite[Lemma 3.2]{tyler2026asymptotics}}]\label{lemma-2.0}
		%		For each $y\ge\frac{\sqrt{3}}{2}$, there exists a $u$ with $|u|<1$ such that
		%		\begin{align*}
			%			|\eta(z)|=\exp\left(-\frac{\pi y}{12}+u 1.01e^{-2\pi y}\right).   
			%		\end{align*}
		%	\end{lemma}
	%	To estimate the Dedekind eta function when $y$ is large, we apply the following lemma.
	%	\begin{lemma}[{\cite[Lemma 2.3]{barman2025lower}}]\label{lemma-2.1}
		%		For every $0<y<\frac{\sqrt{3}}{2}$, one can find a number $v$ satisfying $1<v<1.00873$ such that
		%		\begin{equation*}
			%			\eta(iy) = \exp\left(-\frac{\pi y}{12} - v e^{-2\pi y}\right).
			%		\end{equation*}
		%	\end{lemma}
	For small values of $y > 0$, we apply the functional equation for the Dedekind eta function in the following form.
	
	\begin{lemma}[{\cite[Lemma 3.4]{barman2025lower}}]\label{lemma-2.2}
		Let $0<y<\frac{\sqrt{3}}{2}$, and for each $y$, there exist a $v$ satisfying $1<v<1.01$ such that
		\begin{align*}
			\eta(iy) = y^{-\frac{1}{2}} \exp\left( -\frac{\pi}{12y} - v e^{-\frac{2\pi}{y}} \right).
		\end{align*}
	\end{lemma}	
	%	Note that the group $\text{SL}_{2}(\mathbb{Z})$ acts on the upper half–plane 
	%	$\mathbb{H}=\{\,z \mid \Im z>0\,\}$ through the usual Möbius 
	%	transformations, given by
	%	\begin{equation*}
		%		g z=\frac{az+b}{cz+d} \quad \text{for}\quad g= \begin{pmatrix}
			%			a & b\\
			%			c & d
			%		\end{pmatrix}\in \text{SL}_2(\mathbb{Z}).   
		%	\end{equation*}
	%	The Dedekind eta function satisfies the following identity (see \cite[(3.4)]{tyler2026asymptotics}),
	%	\begin{equation}\label{eq-eta modular}
		%		\begin{aligned}
			%			|\eta(z)|&=\left(\frac{y}{\text{Im}g z}\right)^{-\frac{1}{4}}|\eta(g z)| \quad \text{for all }\quad g\in \text{SL}_{2}(\mathbb{Z}).
			%		\end{aligned}
		%	\end{equation}
	The following lemma helps us to compute the integral in (\ref{eq-ptN1}) in the region $\frac{y}{3} \le |x| \le \frac{1}{2}$.
	\begin{lemma}\label{proposition-error}
		If $y\le \frac{1}{1000}$, then 
		\begin{align*}
			\int_{\frac{y}{3}\le |x|\le \frac{1}{2}}\left|\frac{a_t(z)}{a_t(iy)}\right|dx\ &\le (1.05)^{t} \exp\left(-\frac{\pi t}{120y}\right).
		\end{align*}
	\end{lemma}
	\begin{proof}
		From the definition of $a_t(z)$,
		\begin{align}\label{eq-atzaty}
			\left| \frac{a_t(z)}{a_t(iy)} \right| = \left| \frac{\eta(iy)}{\eta(z)} \right|^{t}.
		\end{align}
		We now show that the above expression is sufficiently small for $\frac{y}{3}\le |x|\le \frac{1}{2}$ and $y\le \frac{1}{1000}$.
		Let $\gamma= \begin{pmatrix}
			a & b\\
			c & d
		\end{pmatrix}\in \text{SL}_2(\mathbb{Z}) $ and $\Im \gamma z\ge \frac{\sqrt{3}}{2}$.
		Proceeding as in the proof of Proposition~3.1 of \cite{barman2026asymptotic} up to equation~(3.3), we obtain
		\begin{align*}
			\left|\frac{\eta(iy)}{\eta(z)}\right|&\le \left(y\Im \gamma z\right)^{-\frac{1}{4}}\exp\left(-\frac{\pi}{12y}+\frac{\pi}{12y}y\Im \gamma z+\frac{1}{228}\right)\\
			&=:\alpha^{-\frac{1}{4}}\exp\left(\frac{\pi\alpha}{12y}\right)\exp\left(-\frac{\pi}{12y}+\frac{1}{228}\right),
		\end{align*}
		where $\alpha=y\Im \gamma z=\frac{y^{2}}{(cx+d)^{2}+c^{2}y^{2}}$. The function $\alpha^{-\frac{1}{4}}\exp\left(\frac{\pi\alpha}{12y}\right)$ reaches its maximum when $\alpha$ is maximum. Note that  $\alpha<\frac{y^{2}}{c^{2}y^{2}}\le \frac{1}{4}$, for $c\ge 2$. When $c = 1$ and $d = 0$, we have $\alpha=\frac{y^{2}}{x^{2}+y^{2}}\le\frac{9}{10}$, since $|x| \ge \frac{y}{3}$. On the other hand, since $\Im \gamma z \ge \frac{\sqrt{3}}{2}$, it follows that $\alpha \in\left[\frac{\sqrt{3}}{2}y, \frac{9}{10}\right]$. Therefore,
		\begin{align*}
			\left|\frac{\eta(iy)}{\eta(z)}\right|&\le\left(\frac{10}{9}\right)^{\frac{1}{4}}\exp\left(-\frac{\pi}{120y}+\frac{1}{228}\right)  \\
			&\le 1.05 \exp\left(-\frac{\pi}{120y}\right).
		\end{align*}
		Using the above estimate in \eqref{eq-atzaty}, the proof follows immediately.
	\end{proof}
		The following lemma will be used in the proof of Proposition~\ref{theorem-main}.
	\begin{lemma}[{\cite[Lemma 4.1]{tyler2026asymptotics}}]\label{lemma-tyler}
		Suppose $\vartheta>38$, $|\rho|<\frac{2}{25}$, $\xi \in(-1,1)$ and $|\theta|\le 1$. Then,
		\begin{align*}
			\int_{-1/3}^{1/3}\exp\left(2\pi i\rho w\right)\exp\left(-\pi\vartheta w^{2}\left(1+2i\xi w+\theta 3w^{2}\right)\right)dw&=\frac{1}{\sqrt{\vartheta}}\left(1+\theta\frac{3.45}{\vartheta}\right).
		\end{align*}
	\end{lemma}
	In \cite{tyler2026asymptotics}, Tyler derived the following explicit formulas for $\mu_m(z)$ according to whether the imaginary part of $z$ is large ($y\ge 1$) or small ($0<y<1$). Detailed proofs of the following two lemmas are given in \cite{barman2026asymptotic}(See Proposition~3.5 and Proposition~3.6).
	\begin{lemma}\label{eq-(2.1)}
		For large $\Im z$, we use the following formula:
		\begin{align*}
			\notag
			\mu_{m}(z)&=\sum_{n=1}^{\infty}z^{m+1}(2\pi in)^{m-1}\sigma(n)\exp(2\pi inz)- \begin{cases} 
				\frac{z^{2}}{24} & \mbox{if } m=0,1 \\ 
				0 & \mbox{if } m\ge 2 
			\end{cases} \quad\text{and}\\  
			\mu_{m}^{\prime}(z)&=\sum_{n=1}^{\infty}\left((2\pi inz)^{m+1}+(m+1)(2\pi inz)^{m}\right) \frac{\sigma(n)}{2\pi in}\exp(2\pi inz)-\begin{cases} 
				\frac{z}{12} & \mbox{if } m=0,1 \\ 
				0 & \mbox{if } m\ge 2. 
			\end{cases}\\ 
			\notag   
		\end{align*}
	\end{lemma}
	\begin{lemma}\label{eq-mue}
		For small $\Im z$, we have
		\begin{align*}
			\mu_{m}(z)&=\sum_{n=1}^{\infty}P_{m}\left(\frac{2\pi in}{z}\right)\sigma(n)\exp\left(-\frac{2\pi in}{z}\right)+\frac{(-1)^{m}m!}{24}+\frac{z}{4\pi i}\begin{cases} 
				\log(-iz) & \mbox{if } m=0 \\ 
				(-1)^{m-1}(m-1)! & \mbox{if } m\ge 1, 
			\end{cases} \\
			\text{and}\quad &\mu_{m}^{\prime}(z)=\sum_{n=1}^{\infty}Q_{m}\left(\frac{2\pi in}{z}\right)\frac{\sigma(n)}{2\pi in}\exp\left(-\frac{2\pi in}{z}\right)+\frac{1}{4\pi i}\begin{cases} 
				1+\log (-iz) & \mbox{if } m=0 \\ 
				(-1)^{m-1}(m-1)! & \mbox{if } m\ge 1, 
			\end{cases}  
		\end{align*}
	\end{lemma}
	where $P_{m}(s)$ is the degree $m-1$ polynomial defined by the recurrence $P_{0}(s)=s^{-1}$ and
	$P_{m}(s)=(s-m)P_{m-1}(s)-sP^{\prime}_{m-1}(s)$.  For $m=0,1,2,3,4$, explicit values of $P_m$ and $Q_m$ are given in \cite{tyler2026asymptotics} (see (4.18)). 
	%	The derivative $\mu^{\prime}_{k}(z)$ is given by:
	%	\begin{equation}
		%		\begin{aligned}
			%			\mu_{m}^{\prime}(z)=&\sum_{n=1}^{\infty}\left((2\pi inz)^{k+1}+(k+1)(2\pi inz)^{k}\right) \frac{\sigma(n)}{2\pi in}\exp(2\pi inz)-\begin{cases} 
				%				z/12 & \mbox{if } k=0 \\ 
				%				0 & \mbox{if } k\ge 1 
				%			\end{cases}\\
			%			&=\sum_{n=1}^{\infty}Q_{k}\left(\frac{2\pi in}{z}\right)\frac{\sigma(n)}{2\pi in}\exp\left(-\frac{2\pi in}{z}\right)+\frac{1}{4\pi i}\begin{cases} 
				%				1+\log (-iz) & \mbox{if } k=0 \\ 
				%				(-1)^{k-1}(k-1)! & \mbox{if } k\ge 1 
				%			\end{cases}   
			%		\end{aligned}
		%	\end{equation} 
	%	where $Q_{k}(r)=r^{2}(P_{k}(r)-P^{\prime}_{k}(r))$. 
	
	Next, we prove the following lemma, which will help us to bound the error of Proposition \ref{theorem-main}.
	\begin{lemma}\label{lemma-mu234}
		Let $|x|<\frac{y}{3}$ and $0<y\le\frac{1}{10}$. Then
		\begin{align*}
			\left|\frac{\mu_{3}(iy)}{\mu_2(iy)}\right|<6\qquad\text{and}  \qquad\left|\frac{\mu_{4}(z)}{\mu_2(iy)}\right|<36.   
		\end{align*}
	\end{lemma}
	\begin{proof}
		From the explicit formula for $\mu_2$ given in Lemma \ref{eq-mue}, we obtain
		\begin{equation}\label{eq-mu2}
			\begin{aligned}
				\mu_{2}(iy)&=\sum_{n=1}^{\infty}\left(\frac{2\pi n}{y}-2\right)\sigma(n) \exp\left(-\frac{2\pi n}{y}\right)+\frac{1}{12}-\frac{y}{4\pi}\\
				&\ge \frac{1}{12}-\frac{y}{4\pi}\ge \frac{1}{12}-\frac{1}{40\pi}\ge \frac{1}{16}\quad\left(\text{as $y\le\frac{1}{10}$}\right).
			\end{aligned} 
		\end{equation}
		By Lemma~4.3$(ii)$ of \cite{tyler2026asymptotics}, we have $-\frac{1}{4}<\mu_3(iy)<0$. Thus, \eqref{eq-mu2} gives
		\begin{align*}
			\left|\frac{\mu_{3}(iy)}{\mu_2(iy)}\right|<6. 
		\end{align*}
		Again, from Lemma \ref{eq-mue}, we have
		\begin{equation}\label{eq-mu42}
			\mu_{4}(z)=\sum_{n=1}^{\infty}\left(\left(\frac{2\pi in}{z}\right)^{3}-12\left(\frac{2\pi in}{z}\right)^{2}+36\left(\frac{2\pi in}{z}\right)-24\right)\sigma(n)\exp\left(-\frac{2\pi in}{z}\right)+1-\frac{3z}{2\pi i}.   
		\end{equation}
		For $r\in \mathbb{Z^{+}}$, we have
		\begin{align*}
			& \left|\left(\frac{2\pi in}{z}\right)^{r}\right| =\frac{(2\pi n)^{r}}{(x^{2}+y^{2})^{r/2}}<\frac{(2\pi n)^{r}}{y^{r}} \quad \text{and}\\
			& \left|\exp\left(-\frac{2\pi in}{z}\right)\right|= \left|\exp\left(-\frac{2\pi in(x-iy)}{(x^{2}+y^{2})}\right)\right|=\exp\left(-\frac{2\pi ny}{(x^{2}+y^{2})}\right).
		\end{align*}
		Since $|x|<\frac{y}{3}$, $\left|\exp\left(-\frac{2\pi in}{z}\right)\right|\le \exp\left(-\frac{9\pi n}{5y}\right).$ Again using the bounds $|x|<\frac{y}{3}$ and $0<y\le \frac{1}{10}$, it follows from \eqref{eq-mu42} and the above equation that
		\begin{align*}%\label{eq-mu45}
			\notag
			|\mu_{4}(z)|
			&\le\sum_{n=1}^{\infty}\left(\frac{(2\pi n)^{3}}{y^{3}}+12\frac{(2\pi n)^{2}}{y^{2}}+36\frac{(2\pi n)}{y}+24\right)\sigma(n)\exp\left(-\frac{9\pi n}{5y}\right)+1+\frac{3y}{2\pi}\sqrt{\frac{10}{9}}\\
			&\le\sum_{n=1}^{\infty}\left((20\pi n)^{3}+12(20\pi n)^{2}+36(20\pi n)+24\right)n^{2}\exp\left(-18\pi n\right)+1+\frac{3}{2\pi}\sqrt{\frac{1}{90}}<2.
		\end{align*}
		Combining (\ref{eq-mu2}), we obtain
		\begin{align*}
			\left|\frac{\mu_{4}(z)}{\mu_2(iy)}\right|<36.     
		\end{align*}
		This completes the proof.
	\end{proof}
	Recall $M=N-\frac{t}{24}$ and let $\theta \in \mathbb{C}$ such that $|\theta|\le 1$. Now, we prove an asymptotic formula for $p_t(N)$. In the proof, the value of $\theta$ may depend on the relevant parameters and may vary from one occurrence to another. In Theorem \ref{theorem-1.2}, we will show that $y$ satisfies $\frac{t\mu_{1}(iy)}{y^{2}}+M=0$. Hence, the restriction \eqref{eq-1.1} on $y$ in the proposition does not affect its generality.
	\begin{proposition}\label{theorem-main} Let $y$ be chosen such that 
		\begin{align}\label{eq-assum}
			\left|\frac{t\mu_{1}(iy)}{y^{2}}+M\right|<\frac{2}{25y},
		\end{align} and assume that $y\le \frac{1}{1000}$. Then
		\begin{align*}
			p_t(N)=\frac{y^{\frac{3}{2}}\exp(2\pi My)a_{t}(iy)}{\sqrt{t\mu_{2}(iy)}}\left(1+\theta\frac{3.48y}{t\mu_2(iy)}\right).    
		\end{align*}
	\end{proposition}
	\begin{proof}
		Recall the integral formula of $p_t(N)$ from (\ref{eq-ptN1}),
		and use Lemma \ref{proposition-error} to write
		\begin{align}\label{eq-p(N,t)}
			p_t(N)&=\exp(2\pi My)a_{t}(iy)\int_{-y/3}^{y/3}\exp\left(-2\pi i Mx+2\pi i\frac{1}{2\pi i}\log\frac{a_{t}(z)}{a_{t}(iy)}\right)dx\\
			\notag
			&+\exp(2\pi My)a_{t}(iy)\left(\theta(1.05)^{t}\exp\left(-\frac{\pi t}{120y}\right)\right).
		\end{align}
		%		Now, applying the Taylor expansion given in (\ref{eq-taylor}), we can express it as follows
		%		\begin{align}\label{eq-taylor err}
			%			\frac{1}{2\pi i}\log \frac{a_{t}(z)}{a_{t}(iy)}
			%			&=tx\frac{\mu_{1}(iy)}{(iy)^{2}}+t\frac{x^{2}}{2}\frac{\mu_{2}(iy)}{(iy)^{3}}\left(1+\frac{x}{3iy}\frac{\mu_{3}(iy)}{\mu_{2}(iy)}+\frac{x^{2}(iy)^{3}}{12(x^{\prime}+iy)^{5}}\frac{\mu_{4}(x^{\prime}+iy)}{\mu_{2}(iy)}\right)
			%		\end{align}
		%		for some $z^{\prime}=x^{\prime}+iy$, where $x^{\prime}$ is between $0$ and $x$. Using Lemma \ref{lemma-mu234}, $\mu_3$ and $\mu_4$ can be bounded by $\mu_2$.
		Recall the Taylor expansion of $\log \frac{a_t(z)}{a_t(iy)}$ from (\ref{eq-taylor}) and apply Lemma \ref{lemma-mu234} to bound $\mu_3$ and $\mu_4$ in term of $\mu_2$ to get  
		\begin{align}\label{eq-tayfinal}
			\frac{1}{2\pi i}\log \frac{a_{t}(z)}{a_{t}(iy)}&=tx\frac{\mu_{1}(iy)}{(iy)^{2}}+t\frac{x^{2}}{2}\frac{\mu_{2}(iy)}{(iy)^{3}}\left(1+2i\xi\frac{x}{y}+3\theta\frac{x^{2}}{y^{2}}\right),   
		\end{align}
		where $\xi \in(-1,1)$ and $|\theta|\le 1.$
		\newline
		Let
		\begin{align}\label{eq-alphaerr}
			\vartheta=\frac{t\mu_{2}(iy)}{y} \qquad\text{and} \qquad\rho=y\left(\frac{t\mu_{1}(iy)}{(iy)^{2}}-M\right). 
		\end{align}
		Now the integrand in the first term of \eqref{eq-p(N,t)} further simplifies after substituting \eqref{eq-tayfinal} to
		\begin{align}\label{eq-exppart}
			\notag
			\exp\left(-2\pi i Mx+2\pi i\frac{1}{2\pi i}\log\frac{a_{t}(z)}{a_{t}(iy)}\right)&=\exp\left(\frac{2\pi ix}{y}y\left(\frac{t\mu_1(iy)}{(iy)^{2}}-M\right)-\pi \frac{x^{2}}{y^{2}}\frac{t\mu_{2}(iy)}{y}\left(1+2i\xi \frac{x}{y}+\theta 3\frac{x^{2}}{y^{2}}\right)\right)\\
			&=\exp\left(\frac{2\pi\rho ix}{y}\right)\exp\left(-\pi\vartheta\frac{x^{2}}{y^{2}}\left(1+2i\xi\frac{x}{y}+\theta 3\frac{x^{2}}{y^{2}}\right)\right).
		\end{align}
		Using the lower bound for $\mu_2$ from (\ref{eq-mu2}), we have
		\begin{align*}
			\vartheta\ge \frac{t}{16y}\ge \frac{1000t}{16}>38 \quad\left(\text{as $y\le 1/1000$ and $t\ge 1$}\right).
		\end{align*}
		Further, by (\ref{eq-assum}), $|\rho|<\frac{2}{25}$. By changing the variable $w=x/y$ in \eqref{eq-exppart} and applying Lemma~\ref{lemma-tyler}, we obtain
		\begin{align*}
			\notag
			\int_{-y/3}^{y/3}\exp\left(-2\pi i Mx+2\pi i\frac{1}{2\pi i}\log\frac{a_{t}(z)}{a_{t}(iy)}\right)dx&=\int_{-1/3}^{1/3}\exp\left(2\pi i\rho w\right)\exp\left(-\pi\vartheta w^{2}\left(1+2i\xi w+\theta 3w^{2}\right)\right)y\,dw\\
			&=\frac{y}{\sqrt{\vartheta}}\left(1+\theta\frac{3.45}{\vartheta}\right).
		\end{align*}
		Combining the above equation with (\ref{eq-p(N,t)}), we obtain 
		\begin{align}\label{eq-pNT}
			p_t(N) &=\exp(2\pi My)a_{t}(iy)\frac{y}{\sqrt{\vartheta}}\left(1+\theta\frac{3.45}{\vartheta}+\frac{\sqrt{\vartheta}}{y}\theta(1.05)^{t}\exp\left(-\frac{\pi t}{120y}\right)\right).
		\end{align}
		Since we assume $\frac{1}{y} \ge 1000$, it follows from (\ref{eq-mu2}) that $\mu_2(iy) < \frac{1}{12}$. Hence, using (\ref{eq-alphaerr}), we derive
		\begin{align*}
			\frac{\vartheta^{\frac{3}{2}}}{y}(1.05)^{t}\exp\left(-\frac{\pi t}{120y}\right)&= \frac{t^{\frac{3}{2}}\mu^{\frac{3}{2}}_2(iy)}{y^{\frac{5}{2}}}(1.05)^{t}\exp\left(-\frac{\pi t}{120y}\right)\\
			&\le \frac{(1000)^{\frac{5}{2}}t^{\frac{3}{2}}}{(12)^{\frac{3}{2}}}(1.05)^{t}\exp\left(-\frac{1000\pi t}{120}\right)<0.03.
		\end{align*}
		%Therefore,
		%	\begin{align*}
			%		\frac{\vartheta^{\frac{1}{2}}}{y}(1.05)^{t}\exp\left(-\frac{\pi t}{120y}\right)&<\frac{0.03}{\vartheta}.
			%	\end{align*}
		Using the above estimate in \eqref{eq-pNT}, we obtain
		\begin{align*}
			p_t(N)&=\exp\left(2\pi My\right)a_t(iy)\frac{y}{\sqrt{\vartheta}}\left(1+\theta\frac{3.48}{\vartheta}\right)
			=\frac{y^{\frac{3}{2}}\exp(2\pi My)a_{t}(iy)}{\sqrt{t\mu_{2}(iy)}}\left(1+\theta\frac{3.48y}{t\mu_2(iy)}\right),    
		\end{align*} 
		as $\vartheta=\frac{t\mu_{2}(iy)}{y}.$
		This completes the proof.
	\end{proof}
	\section{Proof of Theorems \ref{theorem-1.2} and \ref{theorem-1.1}}
	We now turn to the proofs of the main results. 

	\begin{proof}[\textbf{\boldmath Proof of Theorem \ref{theorem-1.2}$(i)$}]
		%		From (\ref{eq-integral}), we have
		%		\begin{align*}%\label{eq-ptN3}
			%			p_{t}(N)&=\int_{-1/2}^{1/2} \exp\left( -2\pi izM\right)a_{t}(z)dx.
			%		\end{align*}
		%		Let
		%		$h(z) = -2\pi i zM + \log a_t(z)
		%		=-2\pi izM-t\log\eta(z).$
		%		As discussed in Section \ref{section1}, the saddle point $y$ is obtained by solving the equation $h'(iy)=0$. Thus,
		To prove the theorem, we determine the saddle point $y$ by solving $\frac{d}{dz}\left(-2\pi iMz+\log a_t(z)\right)=0$ at $z=iy.$ Hence,
		\begin{align*}%\label{eq-saddle}
			2\pi iM+t\frac{d}{dz}\log \eta(z)=0.  
		\end{align*}
		%	By \eqref{eq-muk}, we have
		%	\begin{equation*}
			%		\frac{d}{dz}\log \eta(z)=-\frac{2\pi i}{z^{2}}\mu_{1}(z).   
			%	\end{equation*}
		Note, $\frac{d}{dz}\log \eta(z)=-\frac{2\pi i}{z^{2}}\mu_1(z)$, as defined in~\eqref{eq-muk}, and setting $z=iy$, we obtain
		\begin{equation*}%\label{eq-saddle1}
			\frac{t\mu_{1}(iy)}{y^{2}}=-M=-N+\frac{t}{24}.   
		\end{equation*}
		Next, we prove that the solution $y>0$ is unique. From the explicit expression of $\mu_1$ for large imaginary part as given in Lemma~\ref{eq-(2.1)}, we have 
		\begin{equation*}%\label{eq-pf11}
			\mu_{1}(iy)=\frac{y^{2}}{24}-\sum_{n=1}^{\infty}y^{2}\sigma(n)\exp(-2\pi ny), 
		\end{equation*} 
		and for small imaginary part, Lemma~\ref{eq-mue} gives
		\begin{equation}\label{eq-pf12}
			\mu_{1}(iy)=-\frac{1}{24}+\frac{y}{4\pi}+\sum_{n=1}^{\infty}\sigma(n)\exp\left(-\frac{2\pi n}{y}\right).
		\end{equation}
		From the above two expressions, we obtain 
		\begin{align*}
			\lim_{y\to\infty}\frac{t\mu_1(iy)}{y^{2}}=\frac{t}{24} \quad\text{and}\quad \lim_{y\to 0^{+}}\frac{t\mu_1(iy)}{y^{2}}=-\infty \quad\text{(as $y\le \frac{1}{10}$)}.
		\end{align*}
		Hence, there exists some $y>0$ such that 
		\begin{align*}
			\frac{t\mu_1(iy)}{y^{2}}=-N+\frac{t}{24}.  
		\end{align*}  
		By Lemma~\ref{eq-(2.1)} and Lemma \ref{eq-mue}, the functions $\mu_1$ and $\mu_2$ satisfy the relation $$\mu_2(z)=-2\mu_1(z)+z\mu_1^{\prime}(z).$$ Hence,
		\begin{align*}
			\frac{d}{dy}\left(\frac{t\mu_1(iy)}{y^{2}}\right)=\frac{-t\mu_2(iy)}{y^{3}}<0 \quad(\text{from (\ref{eq-mu2})}).    
		\end{align*}
		Since $\frac{t\mu_1(iy)}{y^{2}}$ is a decreasing function, there exists a unique $y>0$ satisfying~(\ref{eq-1.1}).
		
	\end{proof}

	We now prove the second part of the theorem, namely the asymptotic formula for $p_t(N)$.
	\begin{proof}[\textbf{\boldmath Proof of Theorem \ref{theorem-1.2}~(ii)}]
		We proved that $y$ is a unique solution of $\frac{t\mu_1(iy)}{y^{2}}+M=0$, so it is obvious that $\left|\frac{t\mu_1(iy)}{y^{2}}+M\right|\ll\frac{1}{y}$. From Proposition~\ref{theorem-main}, together with the bound $\mu_2(iy)<\frac{1}{12}$ from \eqref{eq-mu2}, we obtain
		\begin{align*}
			p_t(N)=\frac{y^{\frac{3}{2}}\exp(2\pi My)a_{t}(iy)}{\sqrt{t\mu_{2}(iy)}}\left(1+O\left(\frac{y}{t}\right)\right). 
		\end{align*}  
		%	By (\ref{eq-mu2}), we have $\mu_2(iy)<\frac{1}{12}$, and hence
		%	\begin{align*}
			%		p_t(N)=\frac{y^{\frac{3}{2}}\exp(2\pi My)a_{t}(iy)}{\sqrt{t\mu_{2}(iy)}}\left(1+O\left(\frac{y}{t}\right)\right).   
			%	\end{align*} 
	\end{proof}

\begin{thebibliography}{99}
\bibitem[]{MR644144} James, Gordon; Kerber, Adalbert. The representation theory of the symmetric group. 16 (1981) xxviii+510
\bibitem[]{andrews} Andrews, George E. A survey of multipartitions congruences and identities. surveys in number theory (2008) 1--19
\bibitem[]{atkin} Atkin, AOL. Ramanujan congruences for $p_{-k} (n)$. Canadian Journal of Mathematics 20 (1968) 67--78
\bibitem[]{barman2025lower} Barman, Jayanta; Mahatab, Kamalakshya. Lower Bound for The Number of Zeros in The Character Table of The Symmetric Group. arXiv preprint arXiv:2504.17037 (2025)
\bibitem[]{barman2026asymptotic} Barman, Jayanta; Mahatab, Kamalakshya. Asymptotic Formula for $(t+ 1) $-Regular Partitions. arXiv preprint arXiv:2603.19691 (2026)
\bibitem[]{bouwknegt2002} Bouwknegt, Peter. Multipartitions, generalized {D}urfee squares and affine {L}ie algebra characters. Journal of the Australian Mathematical Society 72 (2002) 395--408
\bibitem[]{bringmann} Bringmann, Kathrin; Kane, Ben; Pahari, Anubhab; Rolen, Larry. Strict log-concavity of k-coloured partitions. Proceedings of the Royal Society of Edinburgh: Section A Mathematics (2026) 1–15
\bibitem[]{bringmann1} Bringmann, Kathrin; Kane, Ben; Rolen, Larry; Tripp, Zack. Fractional partitions and conjectures of {C}hern--{F}u--{T}ang and {H}eim--{N}euhauser. Transactions of the American Mathematical Society, Series B 8 (2021) 615--634
\bibitem[]{chen2014} Chen, William YC; Du, Daniel K; Hou, Qing-Hu; Sun, Lisa H. Congruences of multipartition functions modulo powers of primes. The Ramanujan Journal 35 (2014) 1--19
\bibitem[]{chern} Chern, Shane; Fu, Shishuo; Tang, Dazhao. Some inequalities for k-colored partition functions. The Ramanujan Journal 46 (2018) 713--725
\bibitem[]{fayers2006weights} Fayers, Matthew. Weights of multipartitions and representations of {A}riki--{K}oike algebras. Advances in Mathematics 206 (2006) 112--144
\bibitem[]{gottsche} G{\"o}ttsche, Lothar. The {B}etti numbers of the {H}ilbert scheme of points on a smooth projective surface. Mathematische Annalen 286 (1990) 193--207
\bibitem[]{hardy1918asymptotic} Hardy, Godfrey H; Ramanujan, Srinivasa. Asymptotic formula{\ae} in combinatory analysis. Proceedings of the London Mathematical Society 2 (1918) 75--115
\bibitem[]{meinardus1953} Meinardus, G{\"u}nter. Asymptotische aussagen {\"u}ber partitionen. Mathematische Zeitschrift 59 (1953) 388--398
\bibitem[]{murty2013partition} Murty, M Ram. The partition function revisited. The Legacy of Srinivasa Ramanujan, in: Ramanujan Math. Soc. Lect. Notes Ser., Ramanujan Mathematical Society (2013) 261--279
\bibitem[]{nekrasov} Nekrasov, Nikita A.; Okounkov, Andrei. Seiberg-{W}itten theory and random partitions. the unity of mathematics 244 (2006) 525--596
\bibitem[]{rademacher1938partition} Hans Rademacher. The partition function $p(n)$. Proceedings of the London Mathematical Society 43 (1938) 241--254
\bibitem[]{tyler2026asymptotics} Tyler, Matt. Asymptotics for t-core partitions and {S}tanton's conjecture. Advances in Mathematics 489 (2026) 110805
\end{thebibliography}
\end{document}
